%% kapa: exam=1 round=27 year=2016 period=1 form=A subject=economics subject_ko=경제학원론 seq=29 number=69
%% source: Q-Net 감정평가사 기출문제 (cst00309) artlSeq=5169402 file=extracted/1008/5169402/제27회 감정평가사 1차 문제/제27회 감정평가사 1교시문제(A형).pdf
%% conversion: kordoc 4.24.51 -> tools/gen1.py (auto); review=proofread-sonnet; flags=none
\begin{question}{69}
A국의 총생산함수가 $Y = K^{1/2}L^{1/2}$이다. 이에 관한 설명으로 옳은 것을 모두 고른 것은? (단, $Y$는 국민소득, $K$는 자본량, $L$은 노동량으로 인구와 같다.)
\begin{framedtext}
ㄱ. 총생산함수는 규모에 따른 수익불변의 성질을 가진다. \par
ㄴ. 1인당 자본량이 증가하면 1인당 국민소득은 증가한다. \par
ㄷ. 자본량이 일정할 때, 인구가 증가하면 1인당 국민소득은 감소한다. \par
\end{framedtext}
\begin{choices}
  \item ㄱ
  \item ㄴ
  \item ㄱ, ㄷ
  \item ㄴ, ㄷ
  \item ㄱ, ㄴ, ㄷ
\end{choices}
\end{question}
\answer{5}
